\chapter{Architecture}
\label{chap:architecture}

\section{Invariants}

The project follows these invariants:

\begin{enumerate}
    \item Each piece of application service state is owned by one shard.
    \item Only the owning shard may mutate its service state.
    \item Application services should not be modeled as shared
          \texttt{Arc<Mutex<Service\textgreater{}\textgreater{}} state.
    \item Cross-shard interaction happens through typed messages, typed reply
          handles, or typed submitter calls.
    \item Values crossing shard boundaries are owned.
    \item References to shard-local state must not escape synchronous access
          closures and must not cross \texttt{.await}.
    \item Async work remains on its assigned shard unless explicitly submitted
          elsewhere.
    \item Observability snapshots are owned values, not borrowed runtime
          internals.
    \item Unsafe code is isolated behind small safe APIs.
\end{enumerate}

\section{Workspace layout}

The repository contains four crates:

\begin{description}
    \item[\texttt{sitas-core}] The runtime kernel. Its default build uses
        \texttt{no\_std + alloc} and exposes the portable
        \texttt{ShardRuntime}, reactor, executor, channel, and minimal service
        contracts. Enabling \texttt{std} adds the host runtime, networking,
        observability, and Linux \texttt{io\_uring} modules.
    \item[\texttt{sitas}] The host facade. It enables the \texttt{std} lane,
        re-exports the public API, and owns the examples and integration tests.
    \item[\texttt{sitas-unix}] The Unix implementation of
        \texttt{ShardRuntime}, using threads, a condition-variable parker, and
        the host reactor. It exercises the portable runtime contract on Linux
        and macOS.
    \item[\texttt{sitas-charlotte}] The CharlotteOS backend. It implements the
        same contract with kernel thread and completion-queue syscalls for the
        bare-metal AArch64 target.
\end{description}

The shared contract lives in \filepath{crates/sitas-core/src}. Host examples
live in \filepath{crates/sitas/examples}. Keeping these paths explicit avoids
confusing the portable lane with features available only under \texttt{std}.

\section{The shared-nothing rule}

The central design rule is that application state has one owner. In the
standard-library services that owner is a shard thread receiving typed
commands. In the async runtime it is an executor shard polling futures.

The project therefore avoids
\texttt{Arc<Mutex<Service\textgreater{}\textgreater{}} as its application
model. A mutex serializes mutation but does not identify the responsible shard
or the point where ownership crosses a boundary.

Shared-nothing does not mean synchronization-free. While wakers, reply state,
lifecycle flags, counters, and submission queues use synchronization, application
service state does not. It stays local to its shard.

\begin{notabene}[Why avoid a service mutex?]
A service mutex prevents concurrent mutation. Shard ownership additionally
states which execution context may mutate the value and how other shards
request a change.
\end{notabene}

\begin{center}
\begin{tikzpicture}[
    box/.style={draw=sitasMuted, fill=sitasCodeBackground,
        minimum width=48mm, minimum height=8mm, font=\footnotesize\ttfamily,
        rounded corners=1.5pt},
    arr/.style={->, >=stealth, thick, color=sitasAccent},
    lbl/.style={font=\tiny\itshape, color=sitasMuted}
]
\node[box] (caller) at (0, 0) {Caller thread or task};
\node[box] (shard)  at (0, -1.8) {Owning shard};
\node[box] (reply)  at (0, -3.6) {Owned reply or snapshot};

\draw[arr] (caller) -- node[right, lbl, align=left] {owned command /\\owned future /\\typed submitter call} (shard);
\draw[arr] (shard) -- node[right, lbl] {synchronous local mutation} (reply);
\end{tikzpicture}
\end{center}

\section{Boundaries and lifetimes}

Many runtime bugs occur at boundaries: a transferred value is not owned, a
local reference escapes, a future retains a borrow across \texttt{.await}, or
a completion buffer is freed before the kernel is finished with it.

\texttt{sitas} tries to make those boundaries visible:

\begin{description}
    \item[Cross-shard transfer] uses owned values and explicit
        \texttt{Send + 'static} requirements where a future can move to another
        thread.
    \item[Shard-local access] gives a closure synchronous access to
        \texttt{\&mut T}; the closure returns an owned result and cannot
        suspend.
    \item[Observability] returns owned snapshots. A monitoring thread can keep
        a snapshot after the runtime has moved on.
    \item[Kernel I/O] stores owned buffers until completions are observed or
        until abandoned operation cleanup proves they can be discarded.
\end{description}

\section{Why owned snapshots?}

Observability could expose references into executor internals, but that would
make a diagnostic API participate in the same lifetime and locking problems as
the runtime. The project treats observation as a copy boundary. A snapshot is a
point-in-time value that can be sorted, printed, diffed, or sent to another
thread without keeping the executor borrowed.

Owned snapshots let monitoring code retain and compare observations without
extending the lifetime of task, timer, or shard-local internals.

\section{Layered model}

The runtime can be read as a layered system:

\begin{center}
\begin{tikzpicture}[
    layer/.style={draw=sitasMuted, fill=sitasCodeBackground,
        minimum width=105mm, minimum height=8mm, font=\footnotesize\ttfamily,
        rounded corners=1.5pt},
    arr/.style={->, >=stealth, thick, color=sitasAccent},
    note/.style={font=\tiny\itshape, color=sitasMuted}
]
\node[layer] (api) at (0, 0) {Application service APIs};
\node[layer] (typed) at (0, -1.1) {Typed command/reply and submitter APIs};
\node[layer] (routing) at (0, -2.2) {Placement and shard routing};
\node[layer] (state) at (0, -3.3) {Shard-local service state / ShardLocal<T>};
\node[layer] (executor) at (0, -4.4) {Shard executor or std shard thread};
\node[layer] (runtime) at (0, -5.5) {Mailbox, task queue, timers, readiness, wakeups};
\node[layer] (os) at (0, -6.6) {OS primitives: pipe/poll/epoll/io\_uring/affinity};

\draw[arr] (api) -- (typed);
\draw[arr] (typed) -- (routing);
\draw[arr] (routing) -- (state);
\draw[arr] (state) -- (executor);
\draw[arr] (executor) -- (runtime);
\draw[arr] (runtime) -- (os);

\node[note] at (10mm, -1.65) {owned boundaries};
\node[note] at (10mm, -3.85) {local ownership};
\node[note] at (10mm, -6.05) {runtime mechanics};
\end{tikzpicture}
\end{center}

The standard-library services, host executor, and portable
\texttt{ShardRuntime} lane use different subsets of these layers while sharing
the same ownership rules.

\section{Primary modules}

\begin{description}
    \item[\texttt{shard\_runtime}] The portable contract for shard thread
        creation, typed channels, parking, and per-shard reactors.
        \texttt{ShardSender} is cloneable and the receiver is single-consumer.
        Because the backing ring is single-producer, a small lock serializes
        concurrent senders before they push.
    \item[\texttt{shard\_executor}] The portable single-threaded executor used
        with a \texttt{ReactorBackend}, including the minimal
        \texttt{no\_std} service path.
    \item[\texttt{runtime}] The reusable standard-library kernel for bounded
        shard mailboxes, reply handles, shard startup, and shutdown.
    \item[\texttt{kv}] The reference sharded key-value service.
    \item[\texttt{counter}] A smaller second service proving that the runtime
        primitives are not key-value-specific.
    \item[\texttt{sharded}] A generic \texttt{ShardService} trait and
        \texttt{Sharded<S>} runtime providing common lifecycle (start, shutdown,
        snapshot) for user-defined shard-local services. Shutdown is automatic
        via \texttt{stop\_command}.
    \item[\texttt{placement}] Key-to-shard routing.
    \item[\texttt{stream\_reply}] Streaming reply channels. A
        \texttt{StreamProducer<T>} delivers a sequence of owned values. Two
        async patterns: \texttt{StreamBatch<'a, T>} borrows the reply for
        \texttt{next\_batch()} loops, \texttt{StreamFuture<T>} owns the reply
        with \texttt{collect()} and \texttt{fold()} convenience methods.
        Senders use an atomic refcount so intermediate clone drops do not
        prematurely close the stream.
    \item[\texttt{async\_service}] Adapters that await standard-library
        service replies from executor tasks.
    \item[\texttt{os}] Unix FFI, wake pipes, readiness backends, affinity, and
        Linux \texttt{io\_uring}. Readiness backends cover
        \texttt{epoll} (Linux), \texttt{kqueue} (macOS/iOS/FreeBSD/NetBSD/
        OpenBSD), and \texttt{poll} (other Unix).
    \item[\texttt{executor}] The host single-threaded async kernel, including
        TCP, UDP, spawn backpressure, and Linux \texttt{io\_uring} dispatch.
    \item[\texttt{sharded\_executor}] One executor per shard thread, plus
        submission and observation APIs.
    \item[\texttt{shard\_local}] One owned value per executor shard with
        checked local access. Provides \texttt{with\_current},
        \texttt{with\_current\_result}, \texttt{get\_current},
        \texttt{map\_all}, and worker/scope APIs.
    \item[\texttt{shard\_mailbox}] Typed owned-message transfer between
        shards. Provides uniform key-to-shard routing through
        \texttt{ShardMailboxSet<M>} and non-uniform logical work-unit routing
        through \texttt{WorkUnitMailboxSet<N, M>}.
    \item[\texttt{sharded\_tcp}] Network-facing sharded TCP server with
        \texttt{SO\_REUSEPORT} on Linux, optional Linux
        \texttt{SO\_INCOMING\_CPU} listener placement, an explicit kTLS ULP
        hook for accepted connections, a private socket-options helper for
        \texttt{setsockopt} details, and single-accept distribution on other
        platforms.
    \item[\texttt{metrics}] Thread-safe atomic-counter metrics collector with
        owned \texttt{MetricsSnapshot} for observability.
\end{description}

\section{Portable channel and join safety}

The portable channel is multi-producer and single-consumer. Cloning a sender
does not clone the application value or receiver. Concurrent producers are
serialized around the single-producer ring's push operation; the receiver can
therefore pop without a second queue lock. This synchronization belongs to the
transport boundary, not application service state.

CharlotteOS shard joins use a kernel completion capability registered against
the spawned thread. Each staged closure is keyed by the kernel thread id, so a
thread cannot execute another spawn's closure when scheduling order differs
from creation order. The closure stores its result before exiting; the kernel
then reaps the thread and completes the observer. A join waits for that
completion before taking the result. Generic backends that do not provide a
join still use the explicit always-error \texttt{RawJoinHandle} placeholder.

\section{How to trace a feature through the layers}

When learning or extending the runtime, trace one feature down the stack. For
example, a named task in a scheduling group passes through several layers:

\begin{enumerate}
    \item user code calls a typed spawn API;
    \item the spawner validates that the scheduling group belongs to this
          executor;
    \item a task is allocated with a name and group id;
    \item the scheduler places the task into the group's ready queue;
    \item the executor polls the selected task;
    \item the task records status, wait interest, poll count, and timestamps;
    \item a snapshot later resolves the group id back to an owned group name.
\end{enumerate}

This trace identifies which layer owns validation, scheduling, and
observability, and where the public API crosses a safety boundary.
